\appendix
\section{复现与补充说明}
% Appendix-local numbering: the class restarts the section counter at \appendix
% but not the table counter, which printed "1.1" subsections and a "1.10" table.
\renewcommand{\thesubsection}{\Alph{section}.\arabic{subsection}}
\setcounter{table}{0}
\renewcommand{\thetable}{\Alph{section}\arabic{table}}

\subsection{复现说明}

全部结果由提交材料中的程序生成，运行环境为 Python 3.11.16、NumPy 2.4.6、SciPy 1.17.1、
pandas 3.0.6、PyYAML 6.0.3，所用随机数种子均为固定值。各问的入口程序如下；其中三个复现
入口会完整运行两遍并逐字节比较结果，面板构建程序生成的各表由问题四的复现入口核对不变：

\begin{table}[htbp]
  \centering
  \caption{复现入口}
  \label{tab:repro}
  \begin{tabular}{ll}
    \toprule
    内容 & 入口程序 \\
    \midrule
    问题一：质量评价、冲突与配比 & \texttt{scripts/q1\_reproduce.py} \\
    问题二：标度律与广义形式检验 & \texttt{scripts/q2\_reproduce.py} \\
    问题四：排行榜面板 & \texttt{scripts/q4\_panel\_build.py} \\
    问题四：分解与前沿预测 & \texttt{scripts/q4\_evolution\_reproduce.py} \\
    \bottomrule
  \end{tabular}
\end{table}
% STAGE B: add the Question 3 entry point after T-010 is accepted.

\subsection{人工智能工具使用说明}
% STAGE A: structure only. The content of this subsection must follow the
% organizer's official 2026 regulation on the use of AI tools. That document is
% not in the project's local archive of official material (only the format
% specification and the Word template are), so no rule is inferred from any
% secondary source, and no declaration - of use or of non-use - is written
% here until the official regulation has been obtained and read.
